\section{Stalks}
\label{section-stalks}

\noindent
Let $X$ be a topological space. Let $x \in X$ be a point.
Let $\mathcal{F}$ be a presheaf of sets on $X$.
The {\it stalk of $\mathcal{F}$ at $x$} is the set
$$
\mathcal{F}_x
=
\colim_{x\in U} \mathcal{F}(U)
$$
where the colimit is over the set of open neighbourhoods
$U$ of $x$ in $X$. The set of open neighbourhoods is
partially ordered by (reverse) inclusion:
We say $U \geq U' \Leftrightarrow U \subset U'$.
The transition maps in the system are
given by the restriction maps of $\mathcal{F}$.
See Categories, Section \ref{categories-section-posets-limits}
for notation and terminology regarding (co)limits over systems.
Note that the colimit is a directed colimit.
Thus it is easy to describe $\mathcal{F}_x$. Namely,
$$
\mathcal{F}_x
=
\{
(U, s)
\mid
x\in U, s\in \mathcal{F}(U)
\}/\sim
$$
with equivalence relation given by $(U, s) \sim (U', s')$ if and only if
there exists an open $U'' \subset U \cap U'$ with $x \in U''$ and
$s|_{U''} = s'|_{U''}$. Given a pair $(U, s)$ we sometimes denote
$s_x$ the element of $\mathcal{F}_x$ corresponding to the equivalence
class of $(U, s)$. We sometimes use the phrase
``image of $s$ in $\mathcal{F}_x$'' to denote $s_x$.
For example, given two pairs $(U, s)$ and $(U', s')$ we sometimes
say ``$s$ is equal to $s'$ in $\mathcal{F}_x$'' to indicate
that $s_x = s'_x$. Other authors use the terminology
``germ of $s$ at $x$''.

\medskip\noindent
An obvious consequence of this definition is that
for any open $U \subset X$ there is a canonical map
$$
\mathcal{F}(U)
\longrightarrow
\prod\nolimits_{x \in U} \mathcal{F}_x
$$
defined by $s \mapsto \prod_{x \in U} (U, s)$. Think about it!

\begin{lemma}
\label{lemma-sheaf-subset-stalks}
Let $\mathcal{F}$ be a sheaf of sets on the topological space $X$.
For every open $U \subset X$ the map
$$
\mathcal{F}(U)
\longrightarrow
\prod\nolimits_{x \in U} \mathcal{F}_x
$$
is injective.
\end{lemma}

\begin{proof}
Suppose that $s, s' \in \mathcal{F}(U)$ map to the same element
in every stalk $\mathcal{F}_x$ for all $x \in U$. This means that
for every $x \in U$, there exists an open $V^x \subset U$,
$x \in V^x$ such that $s|_{V^x} = s'|_{V^x}$. But then
$U = \bigcup_{x \in U} V^x$ is an open covering. Thus by the
uniqueness in the sheaf condition we see that $s = s'$.
\end{proof}

\begin{definition}
\label{definition-separated}
Let $X$ be a topological space.
A presheaf of sets $\mathcal{F}$ on $X$ is {\it separated}
if for every open $U \subset X$ the map
$\mathcal{F}(U) \to \prod_{x \in U} \mathcal{F}_x$ is
injective.
\end{definition}

\noindent
Another observation is that the construction of the stalk
$\mathcal{F}_x$ is functorial in the presheaf $\mathcal{F}$.
In other words, it gives a functor
$$
\textit{PSh}(X) \longrightarrow \textit{Sets},
\ \mathcal{F} \longmapsto \mathcal{F}_x.
$$
This functor is called the {\it stalk functor}.
Namely, if $\varphi : \mathcal{F} \to \mathcal{G}$ is
a morphism of presheaves, then we define
$\varphi_x : \mathcal{F}_x \to \mathcal{G}_x$
by the rule $(U, s) \mapsto (U, \varphi(s))$.
To see that this works we have to check that
if $(U, s) = (U', s')$ in $\mathcal{F}_x$ then also
$(U, \varphi(s)) = (U', \varphi(s'))$ in $\mathcal{G}_x$.
This is clear since $\varphi$ is compatible with the
restriction mappings.

\begin{example}
\label{example-stalk-constant-presheaf}
Let $X$ be a topological space. Let $A$ be a set.
Denote temporarily $A_p$ the constant presheaf
with value $A$ ($p$ for presheaf -- not for point).
There is a canonical map of presheaves
$A_p \to \underline{A}$ into the constant sheaf with
value $A$. For every point we have canonical
bijections $A = (A_p)_x = \underline{A}_x$, where
the second map is induced by functoriality from
the map $A_p \to \underline{A}$.
\end{example}



\begin{example}
\label{example-germs-functions}
Suppose $X = \mathbf{R}^n$ with the Euclidean topology.
Consider the presheaf of $\mathcal{C}^\infty$ functions
on $X$, denoted $\mathcal{C}^\infty_{\mathbf{R}^n}$.
In other words, $\mathcal{C}^\infty_{\mathbf{R}^n}(U)$ is the set
of $\mathcal{C}^\infty$-functions $f : U \to \mathbf{R}$.
As in Example \ref{example-basic-continuous-maps}
it is easy to show that this is a sheaf. In fact it
is a sheaf of $\mathbf{R}$-vector spaces.

\medskip\noindent
Next, let $x \in X = \mathbf{R}^n$ be a point. How do
we think of an element in the stalk $\mathcal{C}^\infty_{\mathbf{R}^n, x}$?
Such an element is given by a $\mathcal{C}^\infty$-function
$f$ whose domain contains $x$. And a pair of such
functions $f$, $g$ determine
the same element of the stalk if they agree in a neighbourhood
of $x$. In other words, an element of $\mathcal{C}^\infty_{\mathbf{R}^n, x}$
is the same thing as what is sometimes called
a {\it germ of a $\mathcal{C}^\infty$-function at $x$}.
\end{example}

\begin{example}
\label{example-sheaf-product-pointwise-stalk}
Let $X$ be a topological space. Let $A_x$ be a set for each $x \in X$.
Consider the sheaf $\mathcal{F} : U \mapsto \prod_{x\in U} A_x$ of Example
\ref{example-sheaf-product-pointwise}. We would just like to point
out here that the stalk $\mathcal{F}_x$ of $\mathcal{F}$ at $x$
is in general {\it not} equal to the set $A_x$. Of course there is
a map $\mathcal{F}_x \to A_x$, but that is in general the best you
can say. For example, suppose $x = \lim x_n$ with $x_n \not = x_m$ for
all $n \not = m$ and suppose that $A_y = \{0, 1\}$ for all $y \in X$. Then
$\mathcal{F}_x$ maps onto the (infinite) set of tails of sequences
of $0$s and $1$s. Namely, every open neighbourhood of $x$ contains almost
all of the $x_n$. On the other hand, if every neighbourhood of $x$ contains
a point $y$ such that $A_y = \emptyset$, then $\mathcal{F}_x = \emptyset$.
\end{example}



